% ==========================================
% ADVANCED LATEX TEMPLATE SNIPPETS
% Professional-Grade Code Components
% ==========================================

% ============================================================
% 1. PREMIUM TABLE TEMPLATES
% ============================================================

% --- Triple-Rule Academic Table ---
\begin{table}[H]
\centering
\caption{Academic Standard Table}
\label{tab:academic}
\begin{tabular}{lrrr}
\toprule
\textbf{Variable} & \textbf{Mean} & \textbf{Std. Dev.} & \textbf{N} \\
\midrule
Age & 34.5 & 12.3 & 1,250 \\
Income (\$) & 52,300 & 18,450 & 1,250 \\
Education & 14.2 & 2.8 & 1,250 \\
\bottomrule
\end{tabular}
\end{table}

% --- Gradient Color Table ---
\usepackage{colortbl}
\definecolor{headergrad}{RGB}{68,114,196}
\definecolor{rowalt}{RGB}{242,242,242}

\begin{table}[H]
\centering
\renewcommand{\arraystretch}{1.3}
\begin{tabular}{|l|c|c|c|}
\hline
\rowcolor{headergrad}
\color{white}\textbf{Metric} & \color{white}\textbf{Q1} & \color{white}\textbf{Q2} & \color{white}\textbf{Q3} \\
\hline
\rowcolor{rowalt}
Revenue & \$1.2M & \$1.5M & \$1.8M \\
Profit & \$320K & \$410K & \$520K \\
\rowcolor{rowalt}
Growth & 12\% & 15\% & 18\% \\
\hline
\end{tabular}
\end{table}

% --- Sideways Table for Wide Data ---
\usepackage{rotating}
\begin{sidewaystable}[H]
\centering
\caption{Wide Dataset Overview}
\begin{tabular}{lcccccccc}
\toprule
ID & Param1 & Param2 & Param3 & Param4 & Param5 & Param6 & Param7 & Result \\
\midrule
001 & 23.4 & 45.6 & 78.9 & 12.3 & 34.5 & 56.7 & 89.0 & Pass \\
002 & 34.5 & 56.7 & 89.0 & 23.4 & 45.6 & 67.8 & 90.1 & Pass \\
% ... more rows
\bottomrule
\end{tabular}
\end{sidewaystable}

% --- Table with Conditional Formatting ---
\usepackage{fmtcount}
\newcommand{\highlighthigh}[1]{\textcolor{red}{\textbf{#1}}}
\newcommand{\highlightlow}[1]{\textcolor{green}{#1}}

\begin{table}[H]
\centering
\begin{tabular}{lrr}
\toprule
Region & Sales & Status \\
\midrule
North & \highlighthigh{\$2.5M} & Above Target \\
South & \$1.8M & On Target \\
East & \highlightlow{\$1.2M} & Below Target \\
West & \$2.1M & On Target \\
\bottomrule
\end{tabular}
\end{table}

% ============================================================
% 2. ADVANCED TCOLORBOX COLLECTION
% ============================================================

% --- Learning Objective Box ---
\newtcolorbox{learningobj}[1][]{
  enhanced,
  colback=purple!5!white,
  colframe=purple!75!black,
  coltitle=purple!75!black,
  fonttitle=\bfseries,
  title=Learning Objectives,
  attach boxed title to top left={yshift=-2mm, xshift=2mm},
  boxed title style={colback=purple!75!black},
  before skip=10pt,
  after skip=10pt,
  #1
}

% --- Key Takeaway Box ---
\newtcolorbox{keytakeaway}[1][]{
  enhanced,
  colback=teal!5!white,
  colframe=teal!75!black,
  coltitle=white,
  fonttitle=\bfseries,
  title=Key Takeaway,
  attach boxed title to top left={yshift=-2mm, xshift=2mm},
  boxed title style={colback=teal!75!black},
  before skip=10pt,
  after skip=10pt,
  #1
}

% --- Research Question Box ---
\newtcolorbox{researchquestion}[1][]{
  enhanced,
  colback=magenta!5!white,
  colframe=magenta!75!black,
  coltitle=magenta!75!black,
  fonttitle=\bfseries,
  title=Research Question,
  attach boxed title to top left={yshift=-2mm, xshift=2mm},
  boxed title style={colback=magenta!75!black},
  before skip=10pt,
  after skip=10pt,
  #1
}

% --- Methodology Box ---
\newtcolorbox{methodology}[1][]{
  enhanced,
  colback=brown!5!white,
  colframe=brown!75!black,
  coltitle=brown!75!black,
  fonttitle=\bfseries,
  title=Methodology,
  attach boxed title to top left={yshift=-2mm, xshift=2mm},
  boxed title style={colback=brown!75!black},
  before skip=10pt,
  after skip=10pt,
  #1
}

% --- Proof Box with QED ---
\newtcbtheorem[number within=section]{myproof}{Proof}
  {enhanced jigsaw,
   colback=gray!5!white,
   colframe=gray!75!black,
   coltitle=white,
   fonttitle=\bfseries,
   attach boxed title to top left={yshift=-2mm, xshift=2mm},
   boxed title style={colback=gray!75!black},
   before skip=10pt,
   after skip=10pt}
  {prf}

% --- Quote/Highlight Box ---
\newtcolorbox{quotehighlight}[1][]{
  enhanced,
  colback=orange!5!white,
  colframe=orange!75!black,
  coltitle=orange!75!black,
  fonttitle=\itshape,
  title=Important Quote,
  attach boxed title to top left={yshift=-2mm, xshift=2mm},
  boxed title style={colback=orange!75!black},
  before skip=10pt,
  after skip=10pt,
  #1
}

% ============================================================
% 3. COMPLEX FLOWCHART PATTERNS
% ============================================================

% --- Multi-Branch Decision Tree ---
\begin{tikzpicture}[
    node distance=2cm,
    startstop/.style={rectangle, rounded corners, minimum width=3cm, minimum height=1cm, text centered, draw=black, fill=red!30},
    process/.style={rectangle, minimum width=3cm, minimum height=1cm, text centered, draw=black, fill=blue!20},
    decision/.style={diamond, minimum width=3cm, minimum height=1cm, text centered, draw=black, fill=green!20},
    arrow/.style={thick,->,>=Stealth}
]
    \node [startstop] (start) {Start};
    \node [process, below of=start] (p1) {Process A};
    \node [decision, below of=p1] (d1) {Condition 1?};
    \node [process, below of=d1, yshift=-0.5cm] (p2) {Process B};
    \node [decision, below of=p2, yshift=-0.5cm] (d2) {Condition 2?};
    \node [process, left of=d2, xshift=-2cm] (p3) {Process C};
    \node [process, right of=d2, xshift=2cm] (p4) {Process D};
    \node [startstop, below of=d2, yshift=-1cm] (end) {End};
    
    \path [arrow] (start) -- (p1);
    \path [arrow] (p1) -- (d1);
    \path [arrow] (d1) -- node[right] {Yes} (p2);
    \path [arrow] (d1.east) -- ++(1.5,0) |- (end);
    \path [arrow] (p2) -- (d2);
    \path [arrow] (d2) -- node[above] {Yes} (p4);
    \path [arrow] (d2) -- node[above] {No} (p3);
    \path [arrow] (p3) |- (end);
    \path [arrow] (p4) |- (end);
\end{tikzpicture}

% --- Parallel Process Flow ---
\begin{tikzpicture}[
    box/.style={rectangle, draw, fill=blue!20, text width=5em, text centered, minimum height=2em, rounded corners},
    arrow/.style={thick,->,>=Stealth}
]
    \node [box] (start) {Start};
    \node [box, below of=start, yshift=-1cm] (split) {Split};
    \node [box, left of=split, xshift=-2cm, yshift=-2cm] (path1) {Path 1};
    \node [box, below of=split, yshift=-2cm] (path2) {Path 2};
    \node [box, right of=split, xshift=2cm, yshift=-2cm] (path3) {Path 3};
    \node [box, below of=path2, yshift=-2cm] (merge) {Merge};
    \node [box, below of=merge, yshift=-1cm] (end) {End};
    
    \draw [arrow] (start) -- (split);
    \draw [arrow] (split) -- (path1);
    \draw [arrow] (split) -- (path2);
    \draw [arrow] (split) -- (path3);
    \draw [arrow] (path1) |- (merge);
    \draw [arrow] (path2) -- (merge);
    \draw [arrow] (path3) |- (merge);
    \draw [arrow] (merge) -- (end);
\end{tikzpicture}

% --- V-Model Development Cycle ---
\begin{tikzpicture}[
    phase/.style={rectangle, draw, fill=blue!10, minimum width=3cm, minimum height=1cm, text centered},
    arrow/.style={thick,->,>=Stealth}
]
    % Left side (Requirements)
    \node [phase] (req) {Requirements};
    \node [phase, below of=req] (spec) {System Spec};
    \node [phase, below of=spec] (arch) {Architecture};
    \node [phase, below of=arch] (impl) {Implementation};
    
    % Right side (Testing)
    \node [phase, right of=impl, xshift=4cm] (unit) {Unit Test};
    \node [phase, above of=unit] (int) {Integration};
    \node [phase, above of=int] (sys) {System Test};
    \node [phase, above of=sys] (uat) {UAT};
    
    % Arrows
    \draw [arrow] (req) -- (spec);
    \draw [arrow] (spec) -- (arch);
    \draw [arrow] (arch) -- (impl);
    \draw [arrow] (impl) -- (unit);
    \draw [arrow] (unit) -- (int);
    \draw [arrow] (int) -- (sys);
    \draw [arrow] (sys) -- (uat);
    
    % Connection arrows
    \draw [arrow, dashed] (req.east) -- (uat.west);
    \draw [arrow, dashed] (spec.east) -- (sys.west);
    \draw [arrow, dashed] (arch.east) -- (int.west);
\end{tikzpicture}

% ============================================================
% 4. ADVANCED SYSTEM ARCHITECTURES
% ============================================================

% --- Layered Architecture ---
\begin{tikzpicture}[
    layer/.style={rectangle, draw, fill=blue!20, minimum width=10cm, minimum height=1.5cm, text centered},
    arrow/.style={thick,->,>=Stealth}
]
    \node [layer] (presentation) {Presentation Layer};
    \node [layer, below of=presentation] (business) {Business Logic Layer};
    \node [layer, below of=business] (data) {Data Access Layer};
    \node [layer, below of=data, fill=green!20] (database) {Database};
    
    \draw [arrow, <->] (presentation) -- (business);
    \draw [arrow, <->] (business) -- (data);
    \draw [arrow, <->] (data) -- (database);
\end{tikzpicture}

% --- Event-Driven Architecture ---
\begin{tikzpicture}[
    component/.style={rectangle, draw, fill=blue!10, minimum width=2.5cm, minimum height=1.5cm, text centered},
    eventbus/.style={rectangle, draw, fill=yellow!30, minimum width=8cm, minimum height=1cm, text centered, rounded corners},
    arrow/.style={thick,->,>=Stealth}
]
    \node [component] (producer1) {Producer 1};
    \node [component, right=of producer1] (producer2) {Producer 2};
    \node [eventbus, below of=producer1, yshift=-1cm, xshift=2cm] (bus) {Event Bus};
    \node [component, below of=bus, yshift=-1cm, xshift=-2cm] (consumer1) {Consumer 1};
    \node [component, below of=bus, yshift=-1cm, xshift=2cm] (consumer2) {Consumer 2};
    
    \draw [arrow] (producer1.south) -- (bus.north);
    \draw [arrow] (producer2.south) -- (bus.north);
    \draw [arrow] (bus.south) -- (consumer1.north);
    \draw [arrow] (bus.south) -- (consumer2.north);
\end{tikzpicture}

% --- Cloud Architecture ---
\begin{tikzpicture}[
    cloud/.style={ellipse, draw, fill=blue!10, minimum width=4cm, minimum height=2cm, text centered},
    service/.style={rectangle, draw, fill=green!10, minimum width=2cm, minimum height=1cm, text centered},
    user/.style={rectangle, draw, fill=orange!20, minimum width=1.5cm, minimum height=1.5cm, text centered, rounded corners},
    arrow/.style={thick,->,>=Stealth}
]
    \node [user] (user) {User};
    \node [cloud, right=of user, xshift=2cm] (cloud) {Cloud Platform};
    \node [service, above of=cloud, yshift=-0.5cm] (compute) {Compute};
    \node [service, below of=cloud, yshift=0.5cm] (storage) {Storage};
    \node [service, right of=cloud, xshift=-0.5cm] (network) {Network};
    
    \draw [arrow] (user) -- (cloud);
    \draw [arrow] (cloud) -- (compute);
    \draw [arrow] (cloud) -- (storage);
    \draw [arrow] (cloud) -- (network);
\end{tikzpicture}

% ============================================================
% 5. DATA VISUALIZATION TEMPLATES
% ============================================================

% --- Stacked Bar Chart ---
\begin{tikzpicture}
\begin{axis}[
    ybar stacked,
    enlargelimits=0.15,
    legend style={at={(0.5,-0.15)}, anchor=north,legend columns=-1},
    ylabel={Percentage},
    symbolic x coords={Q1,Q2,Q3,Q4},
    xtick=data,
    nodes near coords,
]
\addplot coordinates {(Q1,30) (Q2,35) (Q3,40) (Q4,45)};
\addplot coordinates {(Q1,20) (Q2,25) (Q3,22) (Q4,20)};
\addplot coordinates {(Q1,15) (Q2,18) (Q3,20) (Q4,22)};
\legend{Category A, Category B, Category C}
\end{axis}
\end{tikzpicture}

% --- Area Chart ---
\begin{tikzpicture}
\begin{axis}[
    area style,
    xlabel={Time},
    ylabel={Value},
    legend pos=north west,
]
\addplot coordinates {(1,2) (2,4) (3,3) (4,7) (5,5)};
\addlegendentry{Series 1}
\addplot coordinates {(1,1) (2,2) (3,4) (4,3) (5,6)};
\addlegendentry{Series 2}
\end{axis}
\end{tikzpicture}

% --- Scatter Plot with Regression ---
\begin{tikzpicture}
\begin{axis}[
    scatter,
    only marks,
    xlabel={X Variable},
    ylabel={Y Variable},
    grid=major,
]
\addplot coordinates {
    (1,2.1) (2,3.9) (3,6.2) (4,7.8) (5,10.1)
    (1.5,2.8) (2.5,4.5) (3.5,6.8) (4.5,8.2) (5.5,10.5)
};
\addplot[red, thick, domain=0:6] {1.7*x + 0.5};
\end{axis}
\end{tikzpicture}

% --- Heatmap ---
\begin{tikzpicture}
\begin{axis}[
    view={0}{90},
    colormap/hot,
    colorbar,
    xlabel={X Axis},
    ylabel={Y Axis},
]
\addplot3[surf, shader=flat corner] coordinates {
    (0,0,1) (1,0,2) (2,0,3) (3,0,4)
    (0,1,2) (1,1,3) (2,1,4) (3,1,5)
    (0,2,3) (1,2,4) (2,2,5) (3,2,6)
    (0,3,4) (1,3,5) (2,3,6) (3,3,7)
};
\end{axis}
\end{tikzpicture}

% ============================================================
% 6. SPECIALIZED DIAGRAMS
% ============================================================

% --- UML Class Diagram ---
\begin{tikzpicture}[
    class/.style={rectangle, draw, fill=white, minimum width=4cm, minimum height=2cm, text centered},
    arrow/.style={thick,->,>=Stealth}
]
    \node [class] (class1) {
        \textbf{ClassName}\\
        \hline
        - privateField: Type\\
        + publicField: Type\\
        \hline
        + method1(): Return\\
        - method2(): void
    };
    \node [class, right=of class1, xshift=2cm] (class2) {
        \textbf{RelatedClass}\\
        \hline
        - data: Type\\
        \hline
        + getData(): Type
    };
    
    \draw [arrow] (class1) -- node[above] {association} (class2);
\end{tikzpicture}

% --- Sequence Diagram ---
\begin{tikzpicture}[
    participant/.style={rectangle, draw, fill=blue!10, minimum width=2cm, minimum height=0.8cm, text centered},
    arrow/.style={thick,->,>=Stealth},
    dashedarrow/.style={thick,->,>=Stealth, dashed}
]
    \node [participant] (user) {User};
    \node [participant, right=of user, xshift=2cm] (api) {API};
    \node [participant, right=of api, xshift=2cm] (db) {Database};
    
    \draw [arrow] (user.south) -- node[right] {Request} (api.south |- 0,-1);
    \draw [arrow] (api.south |- 0,-1) -- node[right] {Query} (db.south |- 0,-2);
    \draw [dashedarrow] (db.south |- 0,-2) -- node[right] {Response} (api.south |- 0,-3);
    \draw [dashedarrow] (api.south |- 0,-3) -- node[right] {Result} (user.south |- 0,-4);
\end{tikzpicture}

% --- Deployment Diagram ---
\begin{tikzpicture}[
    server/.style={rectangle, draw, fill=gray!20, minimum width=3cm, minimum height=2cm, text centered},
    database/.style={cylinder, shape aspect=0.5, draw, fill=green!10, minimum width=2cm, minimum height=1.5cm, text centered},
    arrow/.style={thick,->,>=Stealth}
]
    \node [server] (web) {Web Server};
    \node [server, right=of web, xshift=2cm] (app) {App Server};
    \node [database, below of=app, yshift=-0.5cm] (db) {Database};
    
    \draw [arrow] (web) -- (app);
    \draw [arrow] (app) -- (db);
\end{tikzpicture}

% ============================================================
% 7. MATHEMATICAL VISUALIZATIONS
% ============================================================

% --- Function Graph ---
\begin{tikzpicture}
\begin{axis}[
    axis lines=center,
    xlabel=$x$,
    ylabel=$y$,
    domain=-3:3,
    samples=100,
]
\addplot[blue, thick] {x^2};
\addplot[red, thick, dashed] {2*x};
\end{axis}
\end{tikzpicture}

% --- Parametric Curve ---
\begin{tikzpicture}
\begin{axis}[
    axis lines=center,
    xlabel=$x$,
    ylabel=$y$,
    domain=0:2*pi,
    samples=100,
]
\addplot[blue, thick] ({cos(deg(x))}, {sin(deg(x))});
\end{axis}
\end{tikzpicture}

% --- 3D Surface Plot ---
\begin{tikzpicture}
\begin{axis}[
    view={60}{30},
    xlabel=$x$,
    ylabel=$y$,
    zlabel=$z$,
]
\addplot3[surf, domain=-2:2, samples=20] {x^2 + y^2};
\end{axis}
\end{tikzpicture}

% ============================================================
% 8. CUSTOM LIST ENVIRONMENTS
% ============================================================

\usepackage{enumitem}
\setlist[itemize,1]{label=\color{blue}\textbullet}
\setlist[itemize,2]{label=\color{green}\small\textbullet}
\setlist[enumerate,1]{label=\color{orange}\arabic*.}
\setlist[enumerate,2]{label=\color{purple}\alph*.}

% Custom description list
\newenvironment{keypoints}
  {\begin{description}[font=\bfseries\color{blue}]}
  {\end{description}}

% ============================================================
% 9. PROFESSIONAL HEADERS AND FOOTERS
% ============================================================

\usepackage{fancyhdr}
\pagestyle{fancy}
\fancyhf{}
\fancyhead[L]{\leftmark}
\fancyhead[R]{\thepage}
\fancyfoot[C]{\small Document Title | \today}
\renewcommand{\headrulewidth}{0.4pt}
\renewcommand{\footrulewidth}{0.4pt}

% ============================================================
% 10. CROSS-REFERENCING ENHANCEMENTS
% ============================================================

\usepackage[capitalise,noabbrev]{cleveref}
\Crefname{figure}{Figure}{Figures}
\Crefname{table}{Table}{Tables}
\Crefname{equation}{Equation}{Equations}
\Crefname{section}{Section}{Sections}
\Crefname{theorem}{Theorem}{Theorems}

% Usage examples:
% \Cref{fig:example} produces "Figure 1"
% \Cref{tab:data,fig:chart} produces "Table 1 and Figure 2"

% ============================================================
% 11. BIBLIOGRAPHY STYLES
% ============================================================

\usepackage[style=apa, backend=biber]{biblatex}
\addbibresource{references.bib}

% Alternative styles:
% \usepackage[style=ieee, backend=biber]{biblatex}
% \usepackage[style=chicago-authordate, backend=biber]{biblatex}
% \usepackage[style=vancouver, backend=biber]{biblatex}

% ============================================================
% 12. GLOSSARY AND ACRONYMS
% ============================================================

\usepackage[acronym,toc]{glossaries}
\makeglossaries

\newacronym{api}{API}{Application Programming Interface}
\newacronym{ui}{UI}{User Interface}
\newacronym{ux}{UX}{User Experience}

\newglossaryentry{latex}{
    name=LaTeX,
    description={A document preparation system for high-quality typesetting}
}

% Usage: \gls{api}, \glspl{ui}, \gls{latex}

% ============================================================
% 13. CHEMISTRY AND PHYSICS NOTATION
% ============================================================

\usepackage{chemformula}
\usepackage{siunitx}

% Chemistry examples:
% \ch{H2O}, \ch{CO2}, \ch{C6H12O6}
% \ch{2H2 + O2 -> 2H2O}

% Physics units:
% \SI{9.81}{\meter\per\second\squared}
% \SI{100}{\degreeCelsius}
% \qty{1.6e-19}{\coulomb}

% ============================================================
% 14. MUSIC NOTATION
% ============================================================

\usepackage{musixtex}
% Requires special compilation: latex -> musixflx -> latex

% ============================================================
% 15. ELECTRICAL CIRCUIT DIAGRAMS
% ============================================================

\usepackage[american]{circuitikz}

\begin{circuitikz}
    \draw (0,0) to[battery1] (0,4)
          to[R, l=$R_1$] (4,4)
          to[C, l=$C_1$] (4,0)
          to[ground] (0,0);
\end{circuitikz}

% ============================================================
% 16. GEOGRAPHIC MAPS
% ============================================================

\usepackage{pgfplots}
\pgfplotsset{compat=1.18}
\usepgfplotslibrary{geo}

\begin{tikzpicture}
\begin{axis}[
    geo,
    projection=mercator,
]
\addplot+[only marks] coordinates {
    (longitude1 latitude1)
    (longitude2 latitude2)
};
\end{axis}
\end{tikzpicture}

% ============================================================
% 17. QR CODE GENERATION
% ============================================================

\usepackage{qrcode}
% \qr{https://example.com}
% \qr[level=L]{Text content}
% \qr[level=M, size=2cm]{More content}

% ============================================================
% 18. BARCODE GENERATION
% ============================================================

\usepackage{pst-barcode}
% Requires pst-barcode package and appropriate compiler

% ============================================================
% 19. CALENDAR AND DATE FORMATTING
% ============================================================

\usepackage{calendar}
\usepackage{datetime2}

% Custom date formats:
\DTMnewdatestyle{custom}{%
    \renewcommand{\DTMdisplaydate}[4]{##1/##2/##3}%
}

% ============================================================
% 20. CV/RESUME TEMPLATES
% ============================================================

% Section styling for CVs
\newcommand{\cvsection}[1]{%
    \vspace{1em}%
    \noindent\textbf{\Large #1}%
    \vspace{0.5em}%
    \hrule%
    \vspace{1em}%
}

\newcommand{\cventry}[3]{%
    \noindent\textbf{#1} \hfill \textit{#2}\\%
    \textit{#3}\\%
    \vspace{0.5em}%
}

% ============================================================
% END OF ADVANCED TEMPLATE SNIPPETS
% ============================================================
